\chapter{Implementation}\label{chap:chap4}
In this chapter, we will implement MIRVisualScore by building a visualization example from the ground up. With this practical approach, we will gain insights on how the codebase is structured, how it functions and how it is intended to be used. We will cover everything from the required dependencies, obtaining and preparing the required input files, to creating custom visualizations. By working through this example step-by-step, you will become familiar with the framework's core components and workflows, enabling you to create and adapt visualizations for your own applications. 

\newpage
\section{Codebase Structure}
\begin{Verbatim}
MIRVisualScore/
├── README.md
├── requirements.txt
├── agents.md
├── src/
│   └── functions/
│       ├── __init__.py
│       ├── Beat_Layers.py
│       ├── Spectrogram_layer.py
│       ├── visualization_system.py
│       ├── warp_score.py
│       └── Load_Files.py
├── input_files/
│   ├── beat_this_analysis/
│   └── BWV856/
│       ├── BWV856.mei
│       ├── Performance1/
│       ├── Performance2/
│       └── Performance3/
├── examples/
│   ├── BWV856.py
│   ├── BWV856_EXAMPLE1.py
│   ├── BWV856_EXAMPLE2.py
│   ├── BWV856_COMBINE.py
│   └── Turn_txt_into_MAPS.py
├── output/
└── .gitignore
\end{Verbatim}

\section{Dependencies and Setup}
The depandancies and requirements and how to install them are laid out in the README.md of the codebase. Yet, we will go a bit more in depth in this section. MIRVisualScore as of now, runs using python 3.12.3, using other updated python environments might work, but this was the version used during the writing of this thesis. Using other python versions might cause conflicts with the rest of the codebase. For simplicity, we will run the codebase as if we were running it on a Linux machine, for running other operating systems, simply follow along as the README.md specifies. 

We advise running this codebase with a virtual environment to avoid affecting other projects and repositories you might be working on your system.

Other then the python version the rest of the codebase requires the following dependancies:
\begin{itemize}
    \item matplotlib==3.10.9
    \item numpy==2.4.4
    \item librosa==0.11.0
    \item soundfile==0.13.1
    \item torch==2.11.0
    \item beat\_this==1.1.0
\end{itemize}

To set this up, as the README.md explains, first make sure you have python 3.12.3 installed in your system. Then you can simply run the following commands in the terminal:

\begin{verbatim}
# Create virtual environment
python3.12 -m venv venv

# Activate the virtual enviroment
source venv/bin/activate

#Install requirements
pip install --upgrade pip
pip install -r requirements.txt

\end{verbatim}


\section{Required Files}
For generating visualizations with audio and data aligned with the score we will, at least, three files: A audio file in .wav format with the performance recording, the warped score generated by ScoreWarp, the MAPS file, the .mei score and an .npz with BT data.

\subsection{Getting the Warped Score}
The warped score can be generated using the \href{https://iwk-digital.github.io/scorewarp/}{ScoreWarp online demo}.  In it, one needs to upload the score in .mei format and the MAPS file in order to download the score in svg format. 

The .mei file that contains the score in xml format. The easiest way to get a .mei score is to get the score in .musicxml format and then convert it to .mei, this can be done in multiple ways we will describe two: The first being through the \href{https://editor.verovio.org/}{verovio online editor}. The second, by using musescore's software that allows to then convert to .mei. In either way, it is important to get the original score in .musicxml format to be able to then convert it to .mei using the tools we described. 

\subsection{Getting the MAPS file}
The MAPS file, describes the onsets of the performance and links them to each corresponing score elements of the .mei score through their @xml.id . Getting the maps file might be a bit more cumbersome, but we can describe two different ways of obtaining it: 

The first is the one described in \citet{goeblLetsScoreWarpAgain2025a}, using the MIDI-to-MIDI matching algorithm by \citet{nakamuraPERFORMANCEERRORDETECTION}. However,  this method requires , along side with the .mei score, the performance in midi format. Using examples from the asap dataset \citep{peterAutomaticNoteLevelScoretoPerformance2023} will facilitate this process but if a midi performance is not possible to obtain then the maps file needs to be generated independently.

The MAPS file is at it's core a .json file, and for the purposes of the current MIRVisualScore version, it needs to provide the score's and performance onset data in the following manner:

\begin{verbatim}
{
    "obs_mean_onset": 0.853333333,
    "xml_id": ["x1iw1eq1"],
    "obs_num": 1
}
\end{verbatim}

Where: 

\begin{itemize}
    \item ``\textit{obs\_mean\_onset}'' is the time in seconds of a given onset.
    \item ``\textit{xml\_id}'' is the @xml.id of the corresponding onset note in the .mei score. In the case that the onset has more than one note, like a chord hit, this attribute can have more than one @xml.id. 
    \item ``\textit{obs\_num}'' acts as a simple counter for each detected onset. 
\end{itemize} 

MIRVisualScore depends on the MAPS file for two critical purposes: visualizing onset data and enabling score alignment through ScoreWarp, which uses it to generate the warped score svg. While a comprehensive onset determination across the entire performance is advised for optimal results, ScoreWarp can handle incomplete onset data through interpolation. When onsets are missing from the MAPS file, ScoreWarp automatically interpolates between the provided onset points and their corresponding @xml.id elements to align them with the timeline. 

This flexibility allows for manual annotation at various levels of granularity, depending on the desired alignment precision. For a rough alignment, annotating onsets at structural points (such as every downbeat or every few measures) often suffices, since intermediate notes will be interpolated between these anchors. In the extreme, even providing just the first and last onset positions will correctly position those two points on the timeline, with all intermediate content interpolated between them. This approach is functional but not ideal, as interpolation may introduce inaccuracies in the alignment; however, it offers a practical option when complete onset annotation is not feasible.

For the example that we used in the making of our codebase we modified PP (Piano Precision) to obtain this data. PP, through the ``Piano Aligner'' plugin \citep{jiangPIANOPRECISIONADVANCING2025}, provides onsets and @xml.id of the .mei relative to a performance in .wav in a time instance layer. We made a simple modification where, this time instance layer, instead of displaying each onset with it's corresponding score position, ``measure + position\_in\_measure'', it would display each corresponding @xml.id of each detected onset. With this we could then export the time instances as .txt files with  ``\textit{obs\_mean\_onset}+\textit{xml\_id}'', then converting into the MAPS format using ``\textit{TXT\_to\_Maps()}'' function in @src/functions/warp\_score.

\subsection{Providing BT data}
Since contemporary BT algorithms are ML based, they primarily function by generating a beat activation function, which is a continuous curve representing the probability of a beat occurring at any given frame. Visualizing this underlying probability curve, rather than just the discrete beat assessments, allows us to see the algorithm's internal confidence level, giving crucial insights on the inner workings of the model.

In @src/functions/Beat\_Layers, the function ``\textit{Run\_BeatThis()}'' generates an .npz file that contains the required data for this to be later visualized. This particular function creates the .npz file generated through BeatThis. However, the .npz file format can serve as an intermediate for BT visual layers, as long as it is structered like it is described in ``\textit{Run\_BeatThis()}''. The .npz should contain five key components from the BT model:

\begin{table}[h]
\centering
\begin{tabular}{ll}
\hline
\textbf{Field} & \textbf{Purpose} \\
\hline
\texttt{beat\_times} & Time coordinates (seconds) aligned with probability frames \\
\texttt{beat\_probs} & Continuous beat activation function values \\
\texttt{downbeat\_probs} & Continuous downbeat activation function values \\
\texttt{detected\_beats} & Discrete beat positions derived from peak detection \\
\texttt{detected\_downbeats} & Discrete downbeat positions from peak detection \\
\hline
\end{tabular}
\end{table}

The probabilistic fields capture the algorithm's internal confidence across time, while the detected positions represent the final beat assessments. This separation enables visualization of both the continuous confidence landscape and discrete beat estimates.

\subsection{Summarizing Required Files}
We now know the purpose of each required file and how to get it. To summarize all the files that are directly involved in the example we are constructing:

\begin{table}[h]
\centering
\begin{tabular}{ll}
\hline
\textbf{File} & \textbf{Description} \\
\hline
\texttt{.wav} &  File with the audio recording of the performance \\
\texttt{.mei} &  XML Score of the musical piece\\
\texttt{.maps.json} &  Contains data that links onsets with score elements of the .mei\\
\texttt{.svg} &  The score in visual format generated by ScoreWarp (requires .mei and .maps \textit{a priori})\\
\texttt{.npz} &  File with algorithmic BT data\\
\hline
\end{tabular}
\end{table}

\section{Calling Modules, Methods and Functions}
The codebase divides it's functionalities into 4 main categories, where each, has it's own dedicated python script. They are located inside @src/functions and are:
\begin{itemize}
    \item \texttt{\textbf{visualization\_system.py}}: Focused on establishing the visualization methods for rendering the visual elements, and for managing svg outputs. 
    \item \texttt{\textbf{Spectrogram\_layer.py}}: Responsible for rendering spectrogram visualizations of audio.
    \item \texttt{\textbf{Beat\_Layers.py}}: Dedicated for classes and methods that are involved in the visualization of beat data.
    \item \texttt{\textbf{warp\_score.py}}: For methods that play a rule in the score-to-audio alignment between the warp svg score and the other elements that can be displayed along side the score (the ones generated by the other categories). 
\end{itemize}

In \texttt{\textbf{visualization\_system.py}}, the  ``\texttt{class Layer(ABC)}'' defines the abstract base class for all visual components. This class provides a common interface and structure for creating various visualizations, ensuring consistency and reusability across the codebase. The Layer class defines three essential methods that all subclasses must implement: \texttt{load\_data()}, \texttt{draw()}, and \texttt{to\_svg\_group()}. 

\noindent Each of these three methods serves a distinct purpose in the visualization pipeline:

\begin{itemize}
    \item \texttt{load_data(**kwargs)} — Loads and validates all necessary data for the layer. For instance, \texttt{Spectrogram} computes the mel-scale spectrogram from an audio file, \texttt{BeatProbabilityLayer} loads beat activation functions from an .npz file, and \texttt{Onsets_Layer} extracts onset times from a MAPS JSON file. This method returns a boolean indicating success or failure.
    \item \texttt{draw(ax, shared\_data)} — Renders the visualization onto a matplotlib axis object. It accesses a \texttt{shared\_data} dictionary that contains context information (such as audio duration, file paths, and SVG coordinate mappings) that is shared across all layers. This method returns a tuple of (lines, labels) which are collected for the legend.
    \item \texttt{to\_svg\_group(shared\_data)} — Converts the rendered visualization into vector-based SVG markup. This is essential for creating composited visualizations that can be embedded within the score SVG. Not all layers support SVG export; those that do return SVG group markup, while others return \texttt{None}.
\end{itemize}


\section{Creating the Audio Visualizations}
At it's core, this codebase simply combines matplotlib visualizations and arranges them so that they can be visible along side with a score. 



The reason this codebase is so releant on SVG file formar is because of how Verovio and subsquently ScoreWarp generate the score.

\section{Score and Data Alignment}
In this section we'll go in depth on how the score and data visualization alignment functions at it's core. 

\section{Combining Visualizations}

